% Section 10.1, from lectures/L10/appendix/a_encapsulation_and_polymorphism.md. Each of the three
% demos is typeset in full at the end of the subsection that builds it (§10.1.2, §10.1.5 and
% §10.1.6): the opening paragraphs and the list of files from its README.md, then every source file
% not already printed. The READMEs' long "The example demonstrates:" lists are left out; they
% restate this section.

\section{Encapsulation and Polymorphism in C}
\label{c10:sec:encapsulation}\label{c10:app:a}

\subsection{Why hide a struct's fields}
\label{c10:a:1}
\begin{itemize}
  \item Every driver in Chapter~\ref{c9:ch} is fully exposed: \code{gpio_t} and \code{timer_t}'s
    fields are visible to any file that includes their definition.
  \begin{itemize}
    \item For something as simple as \code{gpio_t}, that's mostly harmless.
    \item For \code{timer_t}, it's less so: nothing stops calling code from reaching in and setting
      its elapsed-time field directly, skipping \code{timer_init()} and quietly breaking the
      invariant \code{timer_running()}/\code{timer_reset()} depend on.
  \end{itemize}
  \item An \textbf{opaque struct} fixes this by splitting a driver across two files:
  \begin{itemize}
    \item A header that only forward-declares the struct's name and lists function prototypes.
    \item A \file{.c} file that defines the struct's actual fields.
    \item Callers only ever see a pointer to an incomplete type: they can pass it around and hand
      it to the driver's functions, but they can't see or touch its fields, so they can't bypass
      the functions meant to maintain its invariants.
  \end{itemize}
  \item Splitting a driver across a header and a \file{.c} file is what makes opacity possible:
  \begin{itemize}
    \item \Secref{c9:a:3} already did this split, for a different reason.
    \item It only works when the struct's real definition lives somewhere the caller doesn't (and
      can't) include.
    \item The next section shows a lighter version of the same idea, hiding \emph{functions}
      rather than fields, before \secref{c10:a:3} hides fields too.
  \end{itemize}
\end{itemize}

\subsection{Hiding functions with \code{static}}
\label{c10:a:2}
\begin{itemize}
  \item Chapter~\ref{c9:ch}'s \code{gpio_t} (\secref{c9:a:2}) stays exactly as public here: every
    field still visible.
  \item What changes is how its behavior is reached:
  \begin{itemize}
    \item \code{gpio_write()}, \code{gpio_read()}, \code{gpio_toggle()}, and
      \code{gpio_enable_pci()} (new in \secref{c9:a:5}'s demo) become \code{static}.
    \item At file scope, \code{static} means internal linkage: invisible outside \file{gpio.c}, a
      different meaning than the storage-duration one it has on a local variable (see
      \secref{c1:a:4}).
    \item \code{gpio_t} gains four function pointers so the now-hidden functions are still
      reachable, just through the struct instead of by name.
  \end{itemize}
\end{itemize}

\noindent\file{driver/gpio.h}, \code{gpio_t}'s fields all still public:

\begin{ccode}
/**
 * @file GPIO driver.
 */
#ifndef DRIVER_GPIO_H_
#define DRIVER_GPIO_H_

#include <stdbool.h>
#include <stdint.h>

/**
 * @brief Enumeration of GPIO modes.
 */
typedef enum
{
    GPIO_MODE_INPUT,        ///< Standard GPIO input.
    GPIO_MODE_INPUT_PULLUP, ///< GPIO input with its internal pull-up enabled.
    GPIO_MODE_OUTPUT,       ///< GPIO output.
} gpio_mode_t;

/**
 * @brief GPIO driver structure.
 */
typedef struct gpio
{
    /** Pointer to the data direction register associated with the GPIO. */
    volatile uint8_t* ddrx;

    /** Pointer to the output register associated with the GPIO. */
    volatile uint8_t* portx;

    /** Pointer to the input register associated with the GPIO. */
    volatile uint8_t* pinx;

    /** Pointer to the pin change mask register associated with the GPIO. */
    volatile uint8_t* pcmskx;

    /** Pin change interrupt control register bit (PCIEx) associated with the GPIO's port. */
    uint8_t pcicrx;

    /** Bit position of the GPIO pin within the associated port. */
    uint8_t pin;

    /** Set state of the given GPIO. Only supported for outputs. */
    void (*write)(struct gpio* self, bool state);

    /** Read state of the given GPIO. */
    bool (*read)(const struct gpio* self);

    /** Toggle state of the given GPIO. Only supported for outputs. */
    void (*toggle)(struct gpio* self);

    /** Enable/disable pin change interrupts for the given GPIO. */
    void (*enable_pci)(struct gpio* self, bool enable);
} gpio_t;

/**
 * @brief Initialize GPIO.
 *
 * @param[in] self GPIO instance.
 * @param[in] pin  Digital pin to use. Must be in range [0, 13].
 * @param[in] mode GPIO mode.
 *
 * @return True on success, false on failure.
 */
bool gpio_init(gpio_t* self, uint8_t pin, gpio_mode_t mode);

#endif /* DRIVER_GPIO_H_ */
\end{ccode}

\begin{itemize}
  \item Each function pointer's \code{self} parameter is typed \code{struct gpio*}, not
    \code{gpio_t*}:
  \begin{itemize}
    \item \code{typedef struct gpio { ... } gpio_t;} only introduces the name \code{gpio_t} once
      the whole declaration finishes, at the final \code{;} after the closing brace.
    \item Inside the braces, while the struct is still being defined, \code{gpio_t} doesn't exist
      yet: writing \code{gpio_t* self} there is a compile error
      (\code{unknown type name 'gpio_t'}).
    \item The tag \code{struct gpio}, on the other hand, is in scope from the moment
      \code{struct gpio \{} appears, so a member can refer to a pointer to its own struct type by
      tag, the same trick a linked list's \code{struct node* next} relies on.
    \item Everywhere \emph{outside} the struct, \code{gpio_t} is the normal name to use; only these
      member declarations, written before their own type has a name, need \code{struct gpio}
      instead.
  \end{itemize}
\end{itemize}
Unlike Chapter~\ref{c9:ch}'s \code{gpio_init(self, mode)}, which expected
\code{ddrx}/\code{portx}/\code{pinx}/\code{pin} already set through a struct literal, this
\code{gpio_init(self, pin, mode)} takes an Arduino pin number and works the registers out itself,
following the pin-number rewrite from Exercise~\ref{c9:ex:4}: pins 0--7 map to \code{PORTD}, pins
8--13 map to \code{PORTB}, and an out-of-range pin returns \code{false} without touching any
register. \file{driver/gpio.c}:

\cfile{lectures/L10/function_pointer_demo/source/driver/gpio.c}

\begin{itemize}
  \item Two fields that used to be \code{const}, \code{pin} and \code{pcicrx}, aren't anymore:
  \begin{itemize}
    \item Both used to be known at the call site and handed in through a struct literal before
      \code{gpio_init()} ever ran.
    \item Now \code{gpio_init()} derives them itself from the pin number argument and assigns them
      with \code{self->pin = ...}/\code{self->pcicrx = ...}.
    \item A \code{const} member can only be set once, by an initializer at the point a struct is
      declared, never by assignment afterward, so once the value is only known and set inside
      \code{gpio_init()}, the field can no longer stay \code{const}. Exercise~\ref{c9:ex:4} asks
      this exact question for the non-function-pointer version of this same rewrite.
  \end{itemize}
\end{itemize}

\begin{console}
function_pointer_demo/
├── Makefile
├── main.c
├── include/
│   └── driver/
│       ├── eeprom.h
│       ├── gpio.h
│       ├── serial.h
│       └── watchdog.h
└── source/
    └── driver/
        ├── eeprom.c
        ├── gpio.c
        ├── serial.c
        └── watchdog.c
\end{console}

\begin{itemize}
  \item \file{driver/gpio.h} and \file{driver/gpio.c} are the full driver shown above; the header
    is printed with its full comments at the end of this section.
  \item \file{driver/eeprom.h}/\file{.c}, \file{driver/serial.h}/\file{.c}, and
    \file{driver/watchdog.h}/\file{.c} are the same EEPROM, UART, and watchdog drivers from
    Chapters~\ref{c6:ch}--\ref{c8:ch}, wired in only so \file{main.c} has something to persist
    state to and report over, unrelated to this chapter's material.
\end{itemize}
\file{main.c} toggles LED1 (\code{D9}) whenever BTN1 (\code{D13}) is pressed, debounces the button
with a 300\,ms timer, persists the new LED state to EEPROM, reports it over serial, and kicks the
watchdog every loop iteration:

\cfile{lectures/L10/function_pointer_demo/main.c}

\begin{itemize}
  \item \code{led1.toggle(&led1)} and \code{btn1.enable_pci(&btn1, true)} instead of
    \code{gpio_toggle(}\allowbreak\code{&led1)} and \code{gpio_enable_pci(&btn1, true)}: the same
    call shape
    \secref{c10:a:6}'s vtables use later, just with the pointers sitting directly on \code{gpio_t}
    instead of behind a separate interface struct.
  \item \code{PCINT0_vect} only toggles the LED, a single register write, and sets
    \code{led1_event}; the EEPROM write and serial report both happen from \code{main()}'s loop
    once \code{led1_event} is seen, following the same ISR-minimality discipline as
    \secref{c6:a:6}: keep the ISR to the minimum, let \code{main()}'s loop handle anything that
    can block.
  \item A second ISR, \code{TIMER0_OVF_vect}, re-enables the button's interrupt after the 300\,ms
    debounce window, using the same \code{enable_pci} function pointer \code{main()} used to arm
    it.
  \item Nothing here hides \code{led1}'s or \code{btn1}'s \emph{fields} yet: \code{led1.ddrx},
    \code{led1.write}, all of it, are exactly as reachable as before. \code{static} only hid the
    \emph{functions}; \secref{c10:a:3} hides the fields too.
\end{itemize}
Build and flash the same way as \secref{c9:a:5}: on Linux, run \code{make} in the demo's
directory, \file{lectures/L10/function_pointer_demo} (\code{make flash} to program the board); on
Windows via Microchip Studio, recreate this folder structure inside the project and add
\file{include} as an include directory.

\subsubsection{The function-pointer demo in full}
The demo in \file{lectures/L10/function_pointer_demo} is about function pointers on a still-public
GPIO driver: a \code{gpio_t} driver where the fields stay fully public, but
\code{write}/\code{read}/\code{toggle}/\code{enable_pci} are reached only through function pointers
stored on the struct itself.

This example demonstrates encapsulation through \code{static} internal linkage rather than through
an opaque struct. The fields of \code{gpio_t}, including its function pointers, remain visible and
assignable from any source file that includes \file{driver/gpio.h}. Only the concrete functions
behind those pointers are hidden, because their definitions have internal linkage inside
\file{gpio.c}.

\code{gpio_init()} accepts an Arduino digital pin number and derives the relevant ATmega328P
registers, bit positions, and pin-change interrupt configuration. It also configures the pin and
assigns all function pointers in a single initialization call.

The implementation is written for compilation on an ATmega328P in Microchip Studio.

\paragraph{Files}
\begin{itemize}
  \item \file{driver/gpio.h}:
  \begin{itemize}
    \item Contains the GPIO driver's public interface.
    \item Defines \code{struct gpio} and its \code{gpio_t} typedef directly in the header.
    \item Exposes every field of the GPIO instance to application code.
    \item Stores the GPIO hardware information required by the internal operations.
    \item Stores the pin's bit number within its hardware port.
    \item Stores the pin-change interrupt group bit used with \code{PCICR}.
    \item Includes four function pointers as struct members:
    \begin{itemize}
      \item \code{write()}
      \item \code{read()}
      \item \code{toggle()}
      \item \code{enable_pci()}
    \end{itemize}
    \item Defines the function-pointer signatures used by the driver operations.
    \item Defines the \code{gpio_mode_t} enum used to select the pin's operating mode.
    \item Declares \code{gpio_init()}.
    \item Exposes \code{gpio_init()} as the driver's only ordinary function with external linkage.
    \item Allows callers to declare individual \code{gpio_t} objects or arrays of GPIO objects.
    \item Requires callers to initialize each object before calling any stored function pointer.
    \item Demonstrates that a public struct provides convention-based rather than enforced
      encapsulation.
  \end{itemize}
  \item \file{driver/gpio.c}:
  \begin{itemize}
    \item Contains the GPIO driver's hardware-specific implementation.
    \item Defines \code{static} implementations of:
    \begin{itemize}
      \item \code{gpio_write()}
      \item \code{gpio_read()}
      \item \code{gpio_toggle()}
      \item \code{gpio_enable_pci()}
    \end{itemize}
    \item Gives these concrete functions internal linkage.
    \item Prevents other translation units from calling the implementations directly by name.
    \item Allows the functions to remain reachable indirectly through the pointers stored in a
      successfully initialized \code{gpio_t}.
    \item Implements \code{gpio_init()}.
    \item Validates the supplied Arduino pin number.
    \item Returns \code{false} when the pin number is outside the supported range.
    \item Performs validation before changing any hardware register.
    \item Maps pins 0--7 to the \code{PORTD} register group.
    \item Maps pins 8--13 to the \code{PORTB} register group.
    \item Calculates the pin's hardware bit position.
    \item Assigns the correct \code{DDRx}, \code{PORTx}, and \code{PINx} register pointers.
    \item Assigns the appropriate pin-change mask register.
    \item Assigns the correct pin-change interrupt group bit.
    \item Configures the pin according to the requested \code{gpio_mode_t}.
    \item Establishes the pin's initial output or pull-up state where required.
    \item Assigns all four internal operation functions to the instance's function-pointer fields.
    \item Leaves the caller-owned \code{gpio_t} ready for use after successful initialization.
    \item Centralizes Arduino-pin-to-ATmega-register mapping in one implementation file.
  \end{itemize}
  \item \file{driver/eeprom.h}/\file{driver/eeprom.c}:
  \begin{itemize}
    \item Contains the EEPROM driver introduced in Chapter~\ref{c7:ch}.
    \item Stores the current LED state in non-volatile memory.
    \item Allows the application to restore the LED state after a reset or power cycle.
    \item Is called only from normal program context.
    \item Remains independent of the GPIO driver.
    \item Does not need access to \code{gpio_t} or its function pointers.
  \end{itemize}
  \item \file{driver/serial.h}/\file{driver/serial.c}:
  \begin{itemize}
    \item Contains the UART driver introduced in Chapter~\ref{c6:ch}.
    \item Initializes serial communication with the configured baud rate.
    \item Reports whether the LED is currently on or off.
    \item Performs output from the main loop rather than from an ISR.
    \item Remains independent of the GPIO implementation.
    \item Demonstrates composition of separate drivers at application level.
  \end{itemize}
  \item \file{driver/watchdog.h}/\file{driver/watchdog.c}:
  \begin{itemize}
    \item Contains the watchdog driver introduced in Chapter~\ref{c8:ch}.
    \item Configures the watchdog during application startup.
    \item Provides the operation used to reset the watchdog counter.
    \item Is serviced once during every main-loop iteration.
    \item Helps recover the system if the main loop stops executing normally.
    \item Operates independently of the GPIO, EEPROM, and UART drivers.
  \end{itemize}
  \item \file{main.c}:
  \begin{itemize}
    \item Declares one caller-owned \code{gpio_t} instance for the LED on D9.
    \item Declares one caller-owned \code{gpio_t} instance for the button on D13.
    \item Initializes the watchdog first, then the GPIOs, the debounce timer and the UART.
    \item Calls \code{gpio_init()} with two pin numbers known to be valid, so it doesn't check the
      results; code that takes its pin numbers from elsewhere must.
    \item Restores the previously saved LED state from EEPROM.
    \item Applies the restored state to the LED.
    \item Arms the button's pin-change interrupt through its \code{enable_pci()} function pointer.
    \item Enables global interrupts only after initialization is complete.
    \item Defines the shared event state used by the ISRs and main loop.
    \item Implements \code{ISR(PCINT0_vect)}:
    \begin{itemize}
      \item Disables the button's pin-change interrupt.
      \item Starts or resets the debounce timing state.
      \item Toggles the LED through the GPIO driver's function pointer.
      \item Sets an event flag.
      \item Performs no EEPROM write.
      \item Performs no UART output.
    \end{itemize}
    \item Implements \code{ISR(TIMER0_OVF_vect)}:
    \begin{itemize}
      \item Tracks the elapsed debounce interval.
      \item Re-enables the button interrupt after approximately 300\,ms.
      \item Stops or resets the debounce timing state when re-arming is complete.
    \end{itemize}
    \item Uses the timer interrupt only for debounce timing.
    \item Resets the watchdog once per main-loop iteration.
    \item Checks the event flag from normal program context.
    \item Clears the event flag once the new state has been written to EEPROM.
    \item Reads the LED's new state.
    \item Writes the new state to EEPROM.
    \item Reports the new state through UART.
    \item Keeps potentially slow operations outside the interrupt service routines.
    \item Coordinates several independent drivers without coupling their implementations together.
  \end{itemize}
\end{itemize}
\file{driver/gpio.c} and \file{main.c} are the files printed above. The header, with its full
comments:

\cfile{lectures/L10/function_pointer_demo/include/driver/gpio.h}

\noindent The serial, EEPROM and watchdog drivers are byte for byte the ones printed in
\secref{c6:a:6}, \secref{c7:a:4} and \secref{c8:a:5}. The demo's \file{Makefile} is the one printed
in \secref{c6:sec:build}, with the same three extra source files in its \code{SRC_FILES} as
Chapter~\ref{c9:ch}'s demo (\secref{c9:a:5}).

\subsection{Opaque structs}
\label{c10:a:3}
\begin{itemize}
  \item Starting over from \secref{c9:a:2}/\secref{c9:a:3}'s \code{gpio_t}, not continuing from
    \secref{c10:a:2}'s version above:
  \begin{itemize}
    \item Hiding fields and hiding functions are independent techniques, not two steps of one
      pipeline, so \code{gpio_write()}/\code{gpio_read()}/\code{gpio_toggle()} go back to being
      ordinary global functions instead of struct fields.
    \item Same fields, same \code{gpio_write()}/\code{gpio_read()}/\code{gpio_toggle()}
      signatures as Chapter~\ref{c9:ch}; \secref{c10:a:2}'s
      \code{enable_pci}/\code{pcmskx}/\code{pcicrx} additions don't carry over either. This
      section is about opacity, not interrupts.
  \end{itemize}
  \item What's new relative to Chapter~\ref{c9:ch}:
  \begin{itemize}
    \item The struct itself moves out of \file{gpio.h} and into \file{gpio.c}.
    \item \code{gpio_init()} is replaced by a \code{malloc()}-based \code{gpio_create()}, since
      callers can no longer declare their own \code{gpio_t} and pass its address once the struct's
      size and layout aren't visible to them.
  \end{itemize}
\end{itemize}

\noindent\file{gpio.h}, the only file calling code includes:

\begin{ccode}
/**
 * @brief GPIO driver for ATmega328p.
 */
#ifndef DRIVER_GPIO_H_
#define DRIVER_GPIO_H_

#include <stdbool.h>
#include <stdint.h>

/**
 * @brief Enumeration of GPIO modes.
 */
typedef enum
{
    GPIO_MODE_INPUT,        ///< Standard GPIO input.
    GPIO_MODE_INPUT_PULLUP, ///< GPIO input with its internal pull-up enabled.
    GPIO_MODE_OUTPUT,       ///< GPIO output.
} gpio_mode_t;

/**
 * @brief GPIO driver structure.
 */
typedef struct gpio gpio_t;   // incomplete type: fields are not visible here

/**
 * @brief Allocate and configure a new GPIO.
 *
 * @param[in] ddrx  Pointer to the data direction register associated with the GPIO.
 * @param[in] portx Pointer to the output register associated with the GPIO.
 * @param[in] pinx  Pointer to the input register associated with the GPIO.
 * @param[in] pin   Bit position of the GPIO pin within the associated port.
 * @param[in] mode  GPIO mode.
 *
 * @return Pointer to the new GPIO, or NULL if allocation failed.
 */
gpio_t* gpio_create(volatile uint8_t* ddrx, volatile uint8_t* portx, volatile uint8_t* pinx,
                     uint8_t pin, gpio_mode_t mode);

/**
 * @brief Set state of the given GPIO.
 *
 * @param[in] self  GPIO instance.
 * @param[in] state GPIO state (true = high, false = low).
 *
 * @note This operation is only supported for GPIOs configured as outputs.
 */
void gpio_write(gpio_t* self, bool state);

/**
 * @brief Read state of the given GPIO.
 *
 * @param[in] self GPIO instance.
 *
 * @return GPIO state (true = high, false = low).
 */
bool gpio_read(const gpio_t* self);

/**
 * @brief Toggle state of the given GPIO.
 *
 * @param[in] self GPIO instance.
 */
void gpio_toggle(gpio_t* self);

#endif /* DRIVER_GPIO_H_ */
\end{ccode}

\noindent\file{gpio.c}, the only file that ever sees \code{gpio_t}'s actual fields:

\begin{ccode}
#include <stddef.h>
#include <stdlib.h>

#include "driver/gpio.h"

struct gpio
{
    volatile uint8_t* ddrx;
    volatile uint8_t* portx;
    volatile uint8_t* pinx;
    uint8_t pin;
};

gpio_t* gpio_create(volatile uint8_t* ddrx, volatile uint8_t* portx, volatile uint8_t* pinx,
                     const uint8_t pin, const gpio_mode_t mode)
{
    gpio_t* self = (gpio_t*)(malloc(sizeof(gpio_t)));
    if (NULL == self) { return NULL; } // Allocation failed.

    self->ddrx = ddrx;
    self->portx = portx;
    self->pinx = pinx;
    self->pin = pin;

    switch (mode)
    {
        case GPIO_MODE_INPUT_PULLUP:
            *self->ddrx &= ~(1U << self->pin);
            *self->portx |= (1U << self->pin);
            break;
        case GPIO_MODE_OUTPUT:
            *self->ddrx |= (1U << self->pin);
            *self->portx &= ~(1U << self->pin);
            break;
        default:
            *self->ddrx &= ~(1U << self->pin);
            *self->portx &= ~(1U << self->pin);
            break;
    }
    return self;
}

void gpio_write(gpio_t* self, const bool state)
{
    if (NULL == self) { return; }
    const bool output = (*self->ddrx & (1U << self->pin));
    if (!output) { return; }

    if (state) { *self->portx |= (1U << self->pin); }
    else { *self->portx &= ~(1U << self->pin); }
}

bool gpio_read(const gpio_t* self)
{
    return (NULL != self) ? (*self->pinx & (1U << self->pin)) : false;
}

void gpio_toggle(gpio_t* self)
{
    if (NULL == self) { return; }
    const bool output = (*self->ddrx & (1U << self->pin));
    if (!output) { return; }
    *self->pinx = (1U << self->pin);
}
\end{ccode}

\begin{itemize}
  \item \code{gpio_create()} calls \code{malloc()} once per instance rather than reserving a
    fixed-size static array:
  \begin{itemize}
    \item A pool sized for the worst case, up to 14 possible Arduino pins, reserves that much RAM
      on every board, even ones that only ever create two or three GPIOs; \code{malloc()} only
      spends what a program actually asks for.
    \item This isn't the ``avoid dynamic memory'' advice AVR programs usually follow, that's
      really about avoiding \emph{repeated} allocation and \code{free()} churn at runtime.
      \code{gpio_create()} is called once per driver instance, at startup, and the instance lives
      for the rest of the program; there's no resizing or fragmentation risk to worry about here.
    \item \code{gpio_create()} returning \code{NULL} on allocation failure is the caller's only
      signal that the heap ran out, same as any other \code{malloc()} call.
  \end{itemize}
\end{itemize}

\subsection{Function pointers}
\label{c10:a:4}
\begin{itemize}
  \item \Secref{c1:a:8} covered the syntax: a variable that holds a function's address, callable
    through the variable exactly like calling the function by name.
  \item What's new here is storing one in a struct, which lets a driver call back into code it
    knows nothing about.
\end{itemize}
Pairing one of \secref{c10:a:3}'s \code{gpio_t} inputs with a callback:

\begin{ccode}
/**
 * @brief Button structure.
 */
typedef struct
{
    /** GPIO associated with the button. */
    gpio_t* pin;

    /**
     * @brief Handle press events.
     */
    void (*on_press)(void); // Called from main(), not the ISR; see L06 Appendix A.6.
} button_t;

/**
 * @brief Check for button press event, invoke its callback if one occurred.
 *
 * @param[in] self Button instance.
 * @param[in] edge True if a press was just detected (e.g. from a flag an ISR set).
 */
static void button_check(const button_t* self, const bool edge)
{
    const bool valid = (NULL != self) && (NULL != self->on_press);
    if (valid && edge) { self->on_press(); }
}
\end{ccode}

\noindent\code{button_t}'s code never mentions what \code{on_press} actually does, toggle an LED,
send a UART message, anything with a matching signature can be plugged in without touching
\code{button_check()}.

\subsection{Worked example: an interrupt-driven \code{gpio_t} with event flags}
\label{c10:a:5}
\begin{itemize}
  \item \Secref{c10:a:4}'s \code{button_t} took an edge flag as a parameter, wherever it came from
    was left unspecified. This demo wires the same idea into a real pin-change interrupt.
  \item \Secref{c10:a:3}'s \code{gpio_t} gains:
  \begin{itemize}
    \item A \code{gpio_callback_t cb} field.
    \item \code{pcmskx}/\code{pcicrx} (the port's pin-change mask register and its enable bit in
      \code{PCICR}).
    \item Three new functions that arm, flag, and dispatch the interrupt.
  \end{itemize}
\end{itemize}
\code{gpio_enable_pci()} arms the interrupt for one GPIO:

\begin{ccode}
void gpio_enable_pci(gpio_t* self)
{
    if (NULL == self) { return; }
    SET(PCICR, self->pcicrx);      // enable interrupts on the associated I/O port
    SET(*self->pcmskx, self->pin); // enable interrupts on the pin
    sei();                         // enable interrupts globally
}
\end{ccode}

\begin{itemize}
  \item Following \secref{c6:a:6}'s discipline, none of the three ISRs (one per I/O port) do
    anything beyond flagging that \emph{something} happened.
  \item \code{event_flags} is \code{volatile} because it's written from an ISR and read from
    \code{main()}'s loop, the exact situation \code{volatile} exists for:
\end{itemize}

\begin{ccode}
/** Event flags for each I/O port (set from an ISR, cleared once handled). */
static volatile uint8_t event_flags = 0U;

static inline void flag_event(const uint8_t port)
{
    if (PORT_COUNT > port) { SET(event_flags, port); }
}

ISR(PCINT0_vect) { flag_event(PORT_OFFSET_B); } // port B: pins 8-13
ISR(PCINT1_vect) { flag_event(PORT_OFFSET_C); } // port C: pins 14-19
ISR(PCINT2_vect) { flag_event(PORT_OFFSET_D); } // port D: pins 0-7
\end{ccode}

\noindent Deciding \emph{what} to do about a flagged event, and actually calling the callback,
happens entirely outside the ISR, in code \code{main()}'s loop drives:

\begin{ccode}
bool gpio_event_occurred(const gpio_t* self)
{
    return NULL != self ? READ(event_flags, self->pcicrx) : false;
}

bool gpio_handle_event(gpio_t* self)
{
    bool handled = false;
    if (gpio_event_occurred(self))
    {
        if (NULL != self->cb)
        {
            cli();
            CLEAR(event_flags, self->pcicrx);
            sei();

            self->cb();
            handled = true;
        }
    }
    return handled;
}
\end{ccode}

\noindent\code{CLEAR()} is a read-modify-write of \code{event_flags}, the same shape as the
\code{PORTB ^=} toggle \secref{c3:a:3} warned about: read the byte, clear one bit, write the byte
back. An ISR that flags another port between that read and that write would have its flag silently
overwritten, so \code{gpio_handle_event()} clears the flag with interrupts disabled. And it clears
it \emph{before} calling the callback, not after: an edge that arrives while the callback runs then
sets the flag again and is handled on the next pass, instead of being wiped out along with the one
just handled.

\code{gpio_event_occurred()} reads \code{event_flags} by \code{pcicrx}
(\code{PCIE0}/\code{PCIE1}/\code{PCIE2}), while the ISRs above flag it by
\code{PORT_OFFSET_B}/\code{_C}/\code{_D} (\code{0}/\code{1}/\code{2}); these only line up because
the ATmega328p happens to number \code{PCIE0}--\code{PCIE2} the same way. Worth noticing when
reading unfamiliar register code: two differently-named constants lining up by coincidence of the
datasheet, not by a shared definition.

\begin{console}
encapsulation_demo/
├── Makefile
├── main.c
├── include/
│   └── driver/
│       ├── gpio.h
│       └── utils.h
└── source/
    └── driver/
        └── gpio.c
\end{console}

\begin{itemize}
  \item \file{driver/gpio.h} and \file{driver/gpio.c}, printed at the end of this section, are the
    full driver, including \code{gpio_new()}/\code{gpio_del()}, extended with the callback and
    interrupt-handling pieces shown above.
  \item \file{driver/utils.h}, printed there too, has the
    \code{SET}/\code{CLEAR}/\code{TOGGLE}/\code{READ} register macros used throughout.
\end{itemize}
\file{main.c} configures pin 9 as an LED output and pin 13 as a button input with its pull-up
enabled, registers \code{button_event()} as the button's callback, and lets pin-change interrupts
drive it:

\cfile{lectures/L10/encapsulation_demo/main.c}

\noindent\code{button_event()} itself never runs inside an ISR, it only ever runs from
\code{gpio_handle_event()}, called from \code{main()}'s loop; the ISRs above never touch
\code{led1}, \code{btn1}, or the callback at all.

Build and flash the same way as \secref{c9:a:5}: on Linux, run \code{make} in the demo's
directory, \file{lectures/L10/encapsulation_demo} (\code{make flash} to program the board); on
Windows via Microchip Studio, recreate this folder structure inside the project and add
\file{include} as an include directory.

\subsubsection{The encapsulation demo in full}
The demo in \file{lectures/L10/encapsulation_demo} is interrupt-driven GPIO with event flags: an
opaque \code{gpio_t} driver extended with a stored callback and pin-change interrupt support.

This example demonstrates encapsulation (\code{gpio_t}'s fields are hidden behind
\file{driver/gpio.h}) combined with a function pointer stored in a struct, dispatched safely from
\code{main()}'s loop. The interrupt service routine only records that an event occurred;
\code{gpio_handle_event()} is what actually invokes the callback.

The implementation is written for compilation on an ATmega328P in Microchip Studio.

\paragraph{Files}
\begin{itemize}
  \item \file{driver/gpio.h}:
  \begin{itemize}
    \item Contains the GPIO driver's public interface.
    \item Forward-declares \code{gpio_t} without exposing the definition of \code{struct gpio}.
    \item Keeps register pointers, pin numbers, callback pointers, and other internal state
      private.
    \item Defines the \code{gpio_mode_t} enum used to select the pin's operating mode.
    \item Declares the lifetime-management functions:
    \begin{itemize}
      \item \code{gpio_new()}
      \item \code{gpio_del()}
    \end{itemize}
    \item Declares the basic GPIO operations:
    \begin{itemize}
      \item \code{gpio_write()}
      \item \code{gpio_read()}
      \item \code{gpio_toggle()}
    \end{itemize}
    \item Declares the pin-change interrupt control functions:
    \begin{itemize}
      \item \code{gpio_enable_pci()}
      \item \code{gpio_disable_pci()}
    \end{itemize}
    \item Declares the event-processing functions:
    \begin{itemize}
      \item \code{gpio_event_occurred()}
      \item \code{gpio_handle_event()}
    \end{itemize}
    \item Provides the only interface that application code needs in order to use the driver.
    \item Prevents application code from depending on the driver's internal representation.
  \end{itemize}
  \item \file{driver/gpio.c}:
  \begin{itemize}
    \item Defines the private \code{struct gpio}.
    \item Stores the hardware information needed to operate each GPIO pin.
    \item Stores the callback function associated with an interrupt-driven GPIO instance.
    \item Implements allocation and initialization performed by \code{gpio_new()}.
    \item Implements cleanup performed by \code{gpio_del()}.
    \item Implements GPIO register access for reading, writing, and toggling pins.
    \item Maps each GPIO pin to the appropriate ATmega328P I/O registers.
    \item Configures the corresponding pin-change interrupt mask bit.
    \item Enables or disables the appropriate pin-change interrupt group.
    \item Maintains the shared event flags used to communicate with the main loop.
    \item Implements \code{gpio_event_occurred()} for checking whether an event is pending.
    \item Implements \code{gpio_handle_event()} for clearing the pending event and invoking the
      stored callback.
    \item Contains the pin-change interrupt handlers:
    \begin{itemize}
      \item \code{ISR(PCINT0_vect)}
      \item \code{ISR(PCINT1_vect)}
      \item \code{ISR(PCINT2_vect)}
    \end{itemize}
    \item Keeps the interrupt handlers intentionally small.
    \item Records events in interrupt context without directly calling application code.
    \item Centralizes all ATmega328P-specific GPIO and interrupt details inside the driver
      implementation.
  \end{itemize}
  \item \file{driver/utils.h}:
  \begin{itemize}
    \item Contains function-like macros for common bit operations.
    \item Provides macros for setting bits in a hardware register.
    \item Provides macros for clearing bits in a hardware register.
    \item Provides macros for toggling bits in a hardware register.
    \item Provides macros for reading individual bits from a hardware register.
    \item Makes low-level register operations more readable.
    \item Reduces repetition in the GPIO driver implementation.
  \end{itemize}
  \item \file{main.c}:
  \begin{itemize}
    \item Demonstrates how application code uses the GPIO driver without accessing its internal
      fields.
    \item Creates one GPIO instance for an LED.
    \item Creates another GPIO instance for a push button.
    \item Configures the LED pin as an output.
    \item Configures the button pin as an input.
    \item Registers a callback for the button GPIO.
    \item Enables the button's pin-change interrupt.
    \item Enables global interrupts through \code{gpio_enable_pci()}, once both GPIOs are
      initialized.
    \item Polls for pending GPIO events in the main loop.
    \item Calls \code{gpio_handle_event()} from normal program context.
    \item Toggles the LED from the registered callback when the button event is handled.
    \item Demonstrates deferred interrupt processing rather than executing the complete response
      inside the ISR.
    \item Releases the GPIO objects through the driver's public interface when cleanup is
      required.
  \end{itemize}
\end{itemize}
\file{main.c} is the file printed above. The driver:

\cfile{lectures/L10/encapsulation_demo/include/driver/gpio.h}
\needspace{8\baselineskip}
\cfile{lectures/L10/encapsulation_demo/include/driver/utils.h}
\needspace{8\baselineskip}
\cfile{lectures/L10/encapsulation_demo/source/driver/gpio.c}

\noindent The demo's \file{Makefile} is the one printed in \secref{c6:sec:build}, with
\file{source/driver/gpio.c} in place of \file{source/driver/serial.c} in its \code{SRC_FILES}:
this demo has no serial driver.

\subsection{Vtables: a \code{gpio_interface_t} interface}
\label{c10:a:6}
\begin{itemize}
  \item \Secref{c10:a:3}'s \code{gpio_t} hardcodes real ATmega328P registers into every instance:
  \begin{itemize}
    \item There's no way to run code that uses it without a chip attached.
    \item There's no way to substitute a fake GPIO for host-side testing.
  \end{itemize}
  \item The fix is a struct of function pointers, an \textbf{interface}, with each concrete
    implementation, real hardware, or an in-memory stub, providing its own vtable:
\end{itemize}

\begin{ccode}
/** GPIO interface. */
typedef struct gpio_interface gpio_interface_t;

/**
 * @brief GPIO virtual table.
 */
typedef struct gpio_vtable
{
    void (*del)(gpio_interface_t** self);
    bool (*read)(const gpio_interface_t* self);
    void (*write)(gpio_interface_t* self, bool value);
    void (*toggle)(gpio_interface_t* self);
} gpio_vtable_t;

/**
 * @brief GPIO interface.
 */
typedef struct gpio_interface
{
    const gpio_vtable_t* vptr;   // pointer to the concrete implementation's vtable
} gpio_interface_t;
\end{ccode}

\noindent Unlike \secref{c10:a:3}/\secref{c10:a:4}'s \code{gpio_t}, \code{gpio_interface_t} has no
register pointers, no pin, nothing hardware-specific at all, just a pointer to a table of
functions. Every operation goes through it: \code{self->vptr->toggle(self)} instead of
\code{gpio_toggle(self)}.

A concrete type implements the interface by embedding \code{gpio_interface_t} as its \emph{first}
field:

\begin{ccode}
typedef struct gpio_impl
{
    gpio_interface_t itf;   // must be first: makes gpio_impl_t* safely reinterpretable as
                             // gpio_interface_t*
    volatile uint8_t* ddrx;
    volatile uint8_t* portx;
    volatile uint8_t* pinx;
    uint8_t pin;
    uint8_t id;
} gpio_impl_t;
\end{ccode}

\begin{itemize}
  \item Because \code{itf} is the struct's first member, C guarantees a \code{gpio_impl_t*} and a
    pointer to its first field share the same address, so a \code{gpio_impl_t*} can be cast to
    \code{gpio_interface_t*} and back without adjusting the pointer.
  \item That's the entire trick a C++ compiler automates for single inheritance: an object's
    address \emph{is} its base class's address, and \code{self->vptr} is exactly the hidden vtable
    pointer every polymorphic C++ object carries.
\end{itemize}
Each vtable function casts back to the concrete type, then operates on its own fields exactly like
\secref{c10:a:3}'s \code{gpio_write()}/\code{gpio_toggle()} did:

\begin{ccode}
static inline gpio_impl_t* get_impl(gpio_interface_t* self) { return (gpio_impl_t*)(self); }

static void gpio_toggle(gpio_interface_t* self)
{
    gpio_impl_t* impl = get_impl(self);
    if ((NULL == impl) || !gpio_impl_is_output(impl)) { return; }
    *(impl->pinx) = (1U << impl->pin); // Writing a 1 to PINx toggles PORTx.
}

static const gpio_vtable_t* gpio_vptr_get_instance(void)
{
    static const gpio_vtable_t vtable = {
        .del = gpio_del, .read = gpio_read, .write = gpio_write, .toggle = gpio_toggle,
    };
    return &vtable;   // one shared vtable; every instance's itf.vptr points at this same address
}

gpio_interface_t* gpio_atmega328p_new(const uint8_t pin, const gpio_mode_t mode)
{
    gpio_impl_t* impl = (gpio_impl_t*)(malloc(sizeof(gpio_impl_t)));
    if (NULL == impl) { return NULL; }
    // ... configure impl->ddrx/portx/pinx/pin from pin, same mapping A.2's gpio_init() did ...
    impl->itf.vptr = gpio_vptr_get_instance();
    return (gpio_interface_t*)(impl);
}
\end{ccode}

\begin{itemize}
  \item Like \secref{c10:a:3}'s \code{gpio_create()}, \code{gpio_atmega328p_new()} calls
    \code{malloc()}, one instance per call, never resized, never repeatedly allocated:
  \begin{itemize}
    \item What's actually new here isn't the allocation strategy, it's what can sit behind the
      pointer: \code{gpio_interface_t*} can point at a \code{gpio_impl_t} or a
      \code{gpio_stub_t}, two differently-sized concrete structs. \Secref{c10:a:3} only ever had
      one struct shape (\code{struct gpio}) to allocate.
  \end{itemize}
  \item \code{del()} takes \code{gpio_interface_t**} rather than \code{gpio_interface_t*}:
  \begin{itemize}
    \item It can \code{free()} the instance and null out the caller's own pointer in one call,
      guarding against a second \code{del()} or a use-after-free through that same variable.
    \item That's a different defense than Exercise~\ref{c9:ex:3} considered for
      \code{gpio_free()} (whether the freed instance's \emph{own} fields should be reset), but the
      same underlying worry: a pointer to a GPIO that no longer safely means what it used to.
  \end{itemize}
\end{itemize}
A second implementation, \code{gpio_stub_t}, satisfies the exact same interface without touching
any hardware, a \code{bool state} field stands in for the registers:

\begin{ccode}
typedef struct gpio_stub
{
    gpio_interface_t itf;   // still first
    bool state;
} gpio_stub_t;

static bool gpio_read(const gpio_interface_t* self)
{
    const gpio_stub_t* impl = (const gpio_stub_t*)(self);
    return NULL != impl ? impl->state : false;
}
\end{ccode}

\noindent Calling code never knows the difference: anything holding a \code{gpio_interface_t*} can
call \code{self->vptr->read(self)}, whether \code{self} came from \code{gpio_atmega328p_new()} or
\code{gpio_stub_new()}.

This is the same mechanism higher-level languages build \code{virtual} functions on top of: an
object with virtual methods carries a hidden pointer to a table of function addresses (its vtable),
and calling a virtual method is, under the hood, exactly the \code{self->vptr->toggle(...)} pattern
above, just generated by the compiler instead of written out by hand.

\begin{console}
interface_demo/
├── Makefile
├── main.c
├── include/
│   └── driver/
│       └── gpio/
│           ├── atmega328p.h
│           ├── interface.h
│           └── stub.h
└── source/
    └── driver/
        └── gpio/
            ├── atmega328p.c
            └── stub.c
\end{console}

\begin{itemize}
  \item \file{driver/gpio/interface.h} is the full interface shown above.
  \item \file{driver/gpio/atmega328p.h}/\file{.c} and \file{driver/gpio/stub.h}/\file{.c} are the
    two implementations in full, all printed at the end of this section.
  \begin{itemize}
    \item Including the pin-registry reservation \code{gpio_atmega328p_new()} uses: the same idea
      as Exercise~\ref{c9:ex:5}'s \code{pin_reg}, just widened from a \code{uint16_t} covering pins
      0--13 to a \code{uint32_t} covering pins 0--19, since this demo's pin mapping also reaches
      \code{PORTC}.
  \end{itemize}
\end{itemize}
\file{main.c} toggles a real LED on D8 whenever a \emph{simulated} button, a \code{gpio_stub_t}
whose state flips every 60,000 main-loop iterations instead of being wired to real hardware,
reports a rising edge. Both constructors can fail, the ATmega328P one also when its pin is already
taken, so \code{main()} checks each result before using it:

\cfile{lectures/L10/interface_demo/main.c}

\noindent Neither \code{simulate_btn_event()} nor \code{main()}'s edge-detection logic cares that
\code{led1} is real hardware and \code{btn1} isn't: both are just \code{gpio_interface_t*}. Swap
\code{gpio_stub_new()} for a second \code{gpio_atmega328p_new()} call and a real button, and
nothing else in this file changes.

Build and flash the same way as \secref{c9:a:5}: on Linux, run \code{make} in the demo's
directory, \file{lectures/L10/interface_demo} (\code{make flash} to program the board); on Windows
via Microchip Studio, recreate this folder structure inside the project and add \file{include} as
an include directory.

\subsubsection{The interface demo in full}
The demo in \file{lectures/L10/interface_demo} is about interfaces and vtables in C: a hand-written
interface (\code{gpio_interface_t}) backed by two independent concrete GPIO implementations ---
real ATmega328P hardware and an in-memory stub --- that satisfy the same contract and can be used
interchangeably.

This example demonstrates, in simplified form, how C++ commonly implements \code{virtual} functions
and single inheritance under the hood. The interface object contains only a pointer to a table of
function pointers. Each concrete type embeds that interface object as its first field, allowing a
pointer to the concrete object to be used as a pointer to the interface without changing its
address.

Calling code holds only a \code{gpio_interface_t*} and dispatches operations through
\code{self->vptr->...}. It does not know whether the object controls a physical ATmega328P pin or
stores its state entirely in memory.

The implementation is written for compilation on an ATmega328P in Microchip Studio.

\paragraph{Files}
\begin{itemize}
  \item \file{driver/gpio/interface.h}:
  \begin{itemize}
    \item Contains the implementation-independent GPIO contract.
    \item Defines \code{gpio_mode_t}.
    \item Defines \code{gpio_vtable_t}.
    \item Declares the complete set of operations that a concrete GPIO implementation must
      provide:
    \begin{itemize}
      \item \code{del()}
      \item \code{read()}
      \item \code{write()}
      \item \code{toggle()}
    \end{itemize}
    \item Defines the signatures of all GPIO operation pointers.
    \item Uses \code{gpio_interface_t*} as the common \code{self} parameter type.
    \item Defines \code{gpio_interface_t}.
    \item Stores only \code{const gpio_vtable_t* vptr} in the interface object.
    \item Contains no ATmega328P register pointers.
    \item Contains no pin number or hardware-specific state.
    \item Contains no stub-specific state.
    \item Allows calling code to dispatch operations without knowing the concrete object layout.
    \item Is the only GPIO header required by code that receives an already-created interface
      object.
    \item Serves a role similar to an abstract base class in C++.
    \item Describes behavior without selecting an implementation.
    \item Allows additional GPIO implementations to be added without changing the
      interface-using logic.
  \end{itemize}
  \item \file{driver/gpio/atmega328p.h}/\file{driver/gpio/atmega328p.c}:
  \begin{itemize}
    \item Declare and define \code{gpio_atmega328p_new(uint8_t pin, gpio_mode_t mode)}.
    \item Provide the concrete implementation backed by real ATmega328P GPIO registers.
    \item Return a \code{gpio_interface_t*} rather than exposing the private concrete type.
    \item Keep \code{gpio_impl_t} private to \file{atmega328p.c}.
    \item Define \code{gpio_impl_t} with \code{gpio_interface_t} as its first field.
    \item Extend the embedded interface with hardware-specific state:
    \begin{itemize}
      \item \code{ddrx}
      \item \code{portx}
      \item \code{pinx}
      \item \code{pin}
      \item \code{id}
    \end{itemize}
    \item Store register pointers for the selected hardware port.
    \item Store the bit position used within those registers.
    \item Store an identifier used to release the pin from the registry.
    \item Define a \code{static const gpio_vtable_t} for the ATmega328P implementation.
    \item Assign the ATmega328P vtable to the embedded interface during construction.
    \item Implement private concrete operations for:
    \begin{itemize}
      \item Reading a hardware pin.
      \item Writing a hardware pin.
      \item Toggling a hardware pin.
      \item Deleting the hardware-backed object.
    \end{itemize}
    \item Convert the incoming interface pointer back to \code{gpio_impl_t*} inside each concrete
      operation.
    \item Validate the supplied Arduino pin number.
    \item Map the pin number to \code{PORTB}, \code{PORTC}, or \code{PORTD}.
    \item Select the corresponding \code{DDRx} and \code{PINx} registers.
    \item Calculate the bit position within the selected port.
    \item Configure the pin according to the requested \code{gpio_mode_t}.
    \item Allocate one private \code{gpio_impl_t} object with \code{malloc()}.
    \item Return \code{NULL} if allocation fails.
    \item Maintain a pin registry representing currently owned physical pins.
    \item Check the registry before claiming a pin.
    \item Return \code{NULL} if another live object already owns the requested pin.
    \item Mark the pin as reserved after successful validation and allocation.
    \item Prevent conflicting GPIO instances from silently modifying the same hardware.
    \item Implement concrete deletion through the interface:
    \begin{itemize}
      \item Validate the supplied pointer.
      \item Recover the concrete \code{gpio_impl_t}.
      \item Reset the pin's data-direction and output state.
      \item Release the pin from \code{pin_registry}.
      \item Free the concrete allocation.
      \item Set the caller's interface pointer to \code{NULL}.
    \end{itemize}
    \item Keep register mapping, reservation, and cleanup details hidden from interface-using
      code.
  \end{itemize}
  \item \file{driver/gpio/stub.h}/\file{driver/gpio/stub.c}:
  \begin{itemize}
    \item Declare and define \code{gpio_stub_new(void)}.
    \item Provide a concrete implementation backed only by memory.
    \item Return the same public type, \code{gpio_interface_t*}, as the hardware constructor.
    \item Keep \code{gpio_stub_t} private to \file{stub.c}.
    \item Define \code{gpio_stub_t} with \code{gpio_interface_t} as its first field.
    \item Extend the embedded interface with a \code{bool state}.
    \item Define a separate \code{static const gpio_vtable_t} for the stub implementation.
    \item Assign the stub vtable during construction.
    \item Allocate one \code{gpio_stub_t} object with \code{malloc()}.
    \item Initialize its logical state.
    \item Return \code{NULL} if allocation fails.
    \item Implement \code{write()} by assigning the in-memory state.
    \item Implement \code{read()} by returning the in-memory state.
    \item Implement \code{toggle()} by inverting the state.
    \item Implement \code{del()} by freeing the stub allocation and nulling the caller's pointer.
    \item Perform no hardware register access.
    \item Require no ATmega328P GPIO pin.
    \item Require no external wiring.
    \item Satisfy the same behavioral contract as the real GPIO implementation.
    \item Allow higher-level logic to be exercised in a deterministic software-only environment.
  \end{itemize}
  \item \file{main.c}:
  \begin{itemize}
    \item Uses \code{gpio_interface_t*} for both GPIO objects.
    \item Creates one real LED with \code{gpio_atmega328p_new()}.
    \item Assigns the LED to Arduino digital pin D8.
    \item Configures the LED pin as an output.
    \item Creates one simulated button with \code{gpio_stub_new()}.
    \item Checks both constructor results before using the returned objects.
    \item Does not access the private fields of either concrete implementation.
    \item Does not need to know the size of either concrete struct.
    \item Calls every GPIO operation through \code{self->vptr}.
    \item Implements \code{simulate_btn_event()}.
    \item Toggles the stub button after approximately every 60,000 main-loop iterations.
    \item Uses the simulated change in state as a substitute for a physical button press.
    \item Stores the previous button state.
    \item Reads the current button state through the interface.
    \item Detects a rising edge when:
    \begin{itemize}
      \item The previous state was low.
      \item The current state is high.
    \end{itemize}
    \item Toggles the physical LED when a rising edge is detected.
    \item Performs both the button read and LED toggle through the same interface contract.
    \item Demonstrates one application simultaneously using two different concrete
      implementations.
    \item Contains no special branch for distinguishing the real object from the stub object.
    \item Demonstrates that dynamic dispatch is determined by each object's vtable.
    \item Deletes both objects through their interface-provided \code{del()} operations.
    \item Leaves implementation-specific cleanup to the concrete types.
    \item Demonstrates hardware-independent application logic built on top of a C interface.
  \end{itemize}
\end{itemize}
\file{main.c} is the file printed above. The interface and its two implementations:

\cfile{lectures/L10/interface_demo/include/driver/gpio/interface.h}
\needspace{8\baselineskip}
\cfile{lectures/L10/interface_demo/include/driver/gpio/atmega328p.h}
\needspace{8\baselineskip}
\cfile{lectures/L10/interface_demo/source/driver/gpio/atmega328p.c}
\needspace{8\baselineskip}
\cfile{lectures/L10/interface_demo/include/driver/gpio/stub.h}
\needspace{8\baselineskip}
\cfile{lectures/L10/interface_demo/source/driver/gpio/stub.c}

\noindent The demo's \file{Makefile} is the one printed in \secref{c6:sec:build}, with
\file{source/driver/gpio/atmega328p.c} and \file{source/driver/gpio/stub.c} in place of
\file{source/driver/serial.c} in its \code{SRC_FILES}: this demo has no serial driver either.
